% demo-cjk2.tex —— M9 中文刀 2 端到端验收样张（UTF-8 输入：源文件直接写中文）
%
% 说明：
% - 与 demo-cjk.tex（刀 1，`\char"<码位>` 形式）对照：本样张开启
%   `\utfinputmode=1` 后**源文件直接写中文**，由输入层把 UTF-8 多字节序列
%   合并为单个 21-bit 字符 token（>255 码位默认 catcode letter，XeTeX 惯例），
%   之后走刀 1 的 OTF 度量 → 折行 → DVI → 渲染 cmap 直查通路。
% - 默认 `\utfinputmode=0`（bytes）时 8-bit 逐字节语义原样保留——
%   TRIP/ETRIP 口径零影响（回归锁：ntex-layout/tests/cjk_charcode.rs）。
% - 字体 FandolSong-Regular（GPL，CTAN fonts/fandol）经 find_otf 搜索链定位。
% - 运行（软光栅 + 真字形）：
%     cargo run -p ntex-backend -- demo-cjk2.tex demo-cjk2 200 --glyphs

\utfinputmode=1
\font\zh=FandolSong-Regular at 24pt
\zh
\hsize=300pt
% 本驱动（ntex-backend）不预载 plain，行距内部参数是引擎默认值；
% 显式给定以得到正常行距（与 plain.tex 的 \baselineskip=12pt / \lineskip=1pt 同理）。
\baselineskip=32pt
\lineskip=4pt
\parindent=0pt

中文排版：源文件直写。

一二三四五六七八九十

% 注：CJK 字符 catcode 为 letter（词），TeX 不会在字母间自动断行——
% CJK 断行规则（\XeTeXlinebreaklocale 类钩子/标点挤压）是依赖链 ④，
% 样张以空行分段使每段不超 \hsize。
\end
